\makeabstract
    {
    Bell-CHSH: It is more than just a Game. It is Entanglement!
    }
    {
    Lautaro Ferreira\textsuperscript{1}, Haochuan Mao\textsuperscript{2}
    }
    {
    \textsuperscript{1} Amherst College, Amherst, MA\\
\textsuperscript{2} Department of Chemistry and Biochemistry, Florida State University, FL
    }
    {
    Experimental violations of the Bell-CHSH inequality have shown that classical mechanics cannot explain quantum-mechanical phenomena. These violations occur in entangled quantum states, which are fundamental to the development of qubits for quantum computing. Particularly, spin-correlated radical pairs (SCRPs) have been recently studied as potential platforms for qubits, because they can be initialized in well-defined pure entangled states at room temperature. The purpose of this research is to simulate an experiment that tracks violations of the Bell-CHSH inequality via the expectation value of a correlator operator C, thereby tracking entanglement and its time evolution using SCRPs.
    }
    {
    The experiments were simulated for a previously studied SCRP initialized in a pure singlet state. Firstly, starting from a proposed spin Hamiltonian operator (H) for SCRPs in the high-field-limit approximation, exponentially decaying time-dependent terms were added to the matrix representation of H in the singlet-triplet basis to simulate how C would evolve over time for individual SCRP states. Secondly, probabilistic mixtures were simulated using a density matrix, and time evolution was tracked using the Liouville-von Neumann equation. Finally, experimental data from microwave pulses used to obtain the density matrix under realistic conditions will be simulated using the EasySpin software in MATLAB.
    }
    {
    Individual-state simulations showed that the decay of the off-diagonal elements increased C to the extent that violations of Bell-CHSH were observed after a short period of time, while the decay of the diagonal elements showed the opposite effect. This allowed us to better understand how SCRP parameters affect entanglement. Simulations using the density matrix showed that, for an SCRP that starts in a pure entangled state, C shows a damped-oscillation behavior, in which violations of Bell-CHSH were observed periodically for a short period of time. We hope to measure this trend using an experimental setup with microwave pulses. Thus, we proposed a potential experiment that initializes the SCRP in a pure entangled state and varies the time until the first microwave pulse to track the time evolution of the density matrix. The sequence of microwave pulses will rotate the density matrix and produce an echo that can be measured as a time-domain spectrum. A Fast Fourier Transform can be used to obtain the frequency-domain spectrum, which can be fitted to extract the appropriate elements of the density matrix to compute C.
    }
    {
    Initial simulations have shown that extracting the density matrix via fitting recovers the simulated trend for C, indicating the potential to generate more realistic data under experimentally realistic conditions. Successfully simulating this data and corresponding it with theoretical results would demonstrate the feasibility of measuring entanglement using SCRPs and microwave pulses, which has never been done before.
    }

    \newpage
    